\documentclass[../../main.tex]{subfiles}
\begin{document}
\section{The static obstruction: quadratic chaos and the two-tail configuration}
\label{sec:qcts}

The time-zero version of two-color covariance control is already a strong statement. It clarifies
why the solved thin-shell theorem is not by itself enough for KLS through this route, and the
quantified two-tail configuration at the end of the section is the static obstruction that both
routes must respect: it fixes the covariance weight of Route~B and marks the boundary of what
model verification can establish for Route~A.

Let $X\sim\nu$ be isotropic and log-concave, and put
\[
  Z=XX^T-I.
\]
For a color
\[
  g=\frac{\one_E-p}{\sqrt{pq}},
  \qquad p=\nu(E),\quad q=1-p,
\]
one computes
\begin{equation}\label{eq:color-tensor}
  \E[gZ]=\sqrt{pq}\left(\Sigma_E-\Sigma_F+(q-p)\delta\delta^T\right).
\end{equation}
At exact balance, $p=q=1/2$, this is $\frac12(\Sigma_E-\Sigma_F)$.

\begin{definition}[Quadratic-chaos thin shell]
\label{def:qcts}
For isotropic log-concave $\nu$, define
\[
  \calQ(\nu)=\sup_{\norm M_{\HS}=1}\Var_\nu(X^TMX).
\]
The quadratic-chaos thin-shell assertion is
\begin{equation}\label{eq:QCTS}
  \calQ(\nu)\le C
\end{equation}
with a universal constant $C$.
\end{definition}

\begin{proposition}[Time-zero equivalence, modulo standard regularization]
\label{prop:qcts-equivalence}
Up to universal constants, \eqref{eq:QCTS} is equivalent to the balanced two-color estimate
\begin{equation}\label{eq:balanced-two-color-static}
  pq\norm{\Sigma_E-\Sigma_F}_{\HS}^2\le C
\end{equation}
for all balanced cuts $E$.
\end{proposition}

\begin{proof}
Assume \eqref{eq:QCTS}. For $\norm M_{\HS}=1$,
\[
  \inner{\E[gZ]}{M}
  =\E\left[g\left(X^TMX-\Tr M\right)\right]
  \le \sqrt{\Var(X^TMX)}\le\sqrt C .
\]
Taking the supremum over $M$ and using \eqref{eq:color-tensor} gives the two-color bound at balance;
the slightly unbalanced case is handled by the explicit $(q-p)\delta\delta^T$ term and
$B\preceq A$.

Conversely, fix $M$ with $\norm M_{\HS}=1$ and set $Y=X^TMX-\Tr M$. In the atomless case choose a
median $m$ of $Y$ and let $E=\{Y\ge m\}$; in the general case use the standard randomized threshold
on the level set $\{Y=m\}$ to obtain exact mass $1/2$. The balanced two-color estimate bounds
$\abs{\E[gY]}$ universally. Since $g$ is the sign of $Y-m$ after randomization, this controls
$\E\abs{Y-m}$ up to a universal factor. Reverse Holder inequalities for bounded-degree polynomials
under log-concave measures, due to Carbery--Wright \cite{CarberyWright2001}, give
$\norm{Y-\E Y}_{L^2}\lesssim\norm{Y-\E Y}_{L^1}$, and hence $\Var(Y)\lesssim C$.
\end{proof}

\subsection{What ordinary thin shell gives}

The thin-shell theorem controls the radial quadratic form:
\[
  \Var(\abs X^2)\le Cn,
\]
which is \eqref{eq:QCTS} only for $M=I/\sqrt n$. Applying thin shell to every projection $P$ gives
\[
  \Var(X^TPX)\le C\operatorname{rank}(P).
\]
For a positive semidefinite matrix $M=\int_0^\infty P_s\dd s$, where
$P_s=\one_{\{M\ge s\}}$, Minkowski yields
\[
  \norm{X^TMX-\Tr M}_{L^2}
  \le C\int_0^\infty\sqrt{\operatorname{rank}(P_s)}\dd s.
\]
The right-hand side is the Lorentz $\ell_{2,1}$ norm of the eigenvalue sequence of $M$, bounded by
$C\sqrt{\log n}\norm M_{\HS}$. Thus projection thin-shell estimates naturally give
\begin{equation}\label{eq:qcts-log}
  \Var(X^TMX)\lesssim \log n\,\norm M_{\HS}^2,
\end{equation}
which matches the logarithmic scale rather than KLS.

\subsection{Projection tests cannot remove the logarithm}

The logarithm is not merely an artifact of the integration. There is an abstract positive operator
$T$ on symmetric matrices such that
\[
  \inner{P}{TP}\le C\operatorname{rank}(P)
\]
for every orthogonal projection $P$, while $\norm T_{\op}\simeq\log n$.

Let
\[
  H_n=\sum_{i=1}^n\frac1i,
  \qquad
  N=\diag\left(\frac1{\sqrt{H_n}},\frac1{\sqrt{2H_n}},\ldots,\frac1{\sqrt{nH_n}}\right),
\]
so that $\norm N_{\HS}=1$, and define
\[
  T(M)=H_n\inner{M}{N}N.
\]
Ky Fan's principle gives, for every rank-$r$ projection $P$,
\[
  \inner{P}{N}\le\sum_{i=1}^r\frac1{\sqrt{iH_n}}\le2\sqrt{\frac r{H_n}}.
\]
Therefore $\inner{P}{TP}\le4r$, whereas $\norm T_{\op}=H_n\simeq\log n$. This is not a log-concave
counterexample. It shows that projection tests alone cannot imply quadratic-chaos thin shell.

\subsection{Static lesson}

A proof of KLS through this route cannot rely only on radial information or on projection tests. It
must either prove the full quadratic-chaos estimate in a way that also controls alignment with
$A_t$, or exploit additional dynamic/geometric information attached to the cut $E$.

\subsection{The quantified two-tail obstruction and the weight calibration}

The following sharpens the qualitative two-tail warning above by tracking the excess of the
configuration, not only its covariance contrast. It is the static obstruction consumed in
covariance-weighted form by Route~B (Section~\ref{sec:stein}) and the configuration whose
dynamical occupation Route~A must control (Section~\ref{sec:carleson}).

\begin{proposition}[Quantified two-tail configuration]
\label{prop:two-tail}
Let $\Lambda\ge1$, $\nu_\Lambda=N(0,\diag(\Lambda,1,\dots,1))$ on $\R^n$, and
\[
  E_\Lambda=\bigl\{x:\abs{x_1}\ge a\sqrt\Lambda\bigr\},
  \qquad a=\Phi^{-1}(3/4)\approx0.6745,
\]
so that $\nu_\Lambda(E_\Lambda)=\tfrac12$. Write $\varphi$ for the standard Gaussian density.
Then:
\begin{enumerate}[label=(\roman*),leftmargin=2.2em]
\item $\delta=0$, hence $r=D=0$ and $K=G$;
\item $G=8a\varphi(a)\,\Lambda\,e_1e_1^T$, hence
$\calS_{\nu_\Lambda}(E_\Lambda)/s=16a^2\varphi(a)^2\Lambda^2\approx0.735\,\Lambda^2$;
\item $\nu_\Lambda^+(E_\Lambda)=2\varphi(a)\Lambda^{-1/2}$ and
$e_0(E_\Lambda)\le2\varphi(a)\Lambda^{-1/2}\approx0.636\,\Lambda^{-1/2}\to0$;
\item the relative excess is bounded below uniformly in $\Lambda$ and $n$:
\[
  \frac{e_0(E_\Lambda)}{I_{\nu_\Lambda}(1/2)}\ \ge\ \frac{2\varphi(a)}{\varphi(0)}-1
  \ \approx\ 0.593 .
\]
\end{enumerate}
Consequently, no inequality of the slice-wise form
$\calS_\nu(E)/s\le C_0+C_1r+\beta D+C_2\,e(E)$, with constants independent of the measure, can
hold for all log-concave $\nu$ and all sets of mass $\tfrac12$.
\end{proposition}

\begin{proof}
(i) is symmetry. (ii): conditioning on $\{\abs Z\ge a\}$ for $Z=X_1/\sqrt\Lambda\sim N(0,1)$
affects only the first coordinate, and the Gaussian integrations
$\int_a^\infty z^2\varphi=a\varphi(a)+1-\Phi(a)$, $1-\Phi(a)=\tfrac14$, give
$\E[Z^2\mid\abs Z\ge a]=1+4a\varphi(a)$ and $\E[Z^2\mid\abs Z<a]=1-4a\varphi(a)$, hence
$G_{11}=8a\varphi(a)\Lambda$ with all other entries zero; then
$\calS/s=s\norm G_\HS^2=\tfrac14\cdot64a^2\varphi(a)^2\Lambda^2$. (iii): the boundary consists
of the two hyperplanes $\{x_1=\pm a\sqrt\Lambda\}$, each of boundary density
$\varphi(a)\Lambda^{-1/2}$; and $e_0\le\nu^+$ since $I\ge0$. (iv):
$I_{\nu_\Lambda}(\tfrac12)\le\varphi(0)\Lambda^{-1/2}$ by the halfspace $\{x_1\le0\}$, so
$e_0/I(\tfrac12)=\nu^+/I(\tfrac12)-1\ge2\varphi(a)/\varphi(0)-1$. The final claim follows by
letting $\Lambda\to\infty$: the left side grows like $\Lambda^2$, the right side is
$C_0+C_2O(\Lambda^{-1/2})$.
\end{proof}

\begin{remark}[Weight calibration and the rate $T^{1+\gamma}$]
\label{rem:weight-explains-rate}
Three consequences.

\emph{(a) No slice-wise proof.} At $t=0$ the initial measure is isotropic and the configuration
of Proposition~\ref{prop:two-tail} cannot occur; but anisotropic posteriors
$\mu_t$ with $\lmax(A_t)=\Lambda\gg1$ are exactly the states the dynamics may visit, and there
the slice-wise inequality fails. Any proof of the Stein-trace estimate must therefore show that
such configurations carry negligible expected occupation --- a covariance-occupation statement
of the same nature as the Carleson problem itself. This is the precise sense in which the
excess-propagation question and the operator-to-trace upgrade are two faces of one obstruction.

\emph{(b) The weight $(1+\norm A_\op)^{5/2}$ is calibrated, not chosen.} Matching powers of
$\Lambda$ ($\calS/s\asymp\Lambda^2$ against $e\asymp\Lambda^{-1/2}$) identifies
$(1+\norm{A}_\op)^{5/2}$ as the minimal covariance weight rendering the excess term statically
consistent: indeed $e_0\cdot\Lambda^{5/2}=2\varphi(a)\Lambda^2$ matches
$\calS/s\approx0.735\Lambda^2$ up to the universal factor $8a^2\varphi(a)\approx1.16$. The
uniform lower bound (iv) on the \emph{relative} excess is consistent evidence: two-tail sets
are never relatively near-minimal in the Gaussian model, so a stability theory normalized at
the Cheeger scale of the posterior is not contradicted by the example.

\emph{(c) Why $T^{1+\gamma}$.} By the Brascamp--Lieb cap \eqref{eq:BL-cap} the weight is at most
$(1+t^{-1})^{5/2}$ pathwise. A weighted excess bound of Carleson type can therefore hold near
$t=0$ only if $\E e_t$ decays at a definite rate as $t\downarrow0$ in the dangerous regimes
(heuristically $\E e_t\lesssim t^{3/2}$ where the cap saturates), or if large $\norm{A_t}_\op$
at small times is correspondingly rare. The exponent $T^{1+\gamma}$ in
\eqref{eq:intro-weighted-excess} is the shadow of this weight: it was postulated in the
original formulation without the weight, and the present calibration explains its necessity.
\end{remark}
\end{document}
